\documentclass[12pt, titlepage]{article}

\usepackage{amsmath, mathtools}
\usepackage[round]{natbib}
\usepackage{amsfonts}
\usepackage{amssymb}
\usepackage{graphicx}
\usepackage{colortbl}
\usepackage{xr}
\usepackage{hyperref}
\usepackage{longtable}
\usepackage{xfrac}
\usepackage{tabularx}
\usepackage{float}
\usepackage{siunitx}
\usepackage{booktabs}
\usepackage{multirow}
\usepackage[section]{placeins}
\usepackage{caption}
\usepackage{fullpage}

\hypersetup{
bookmarks=true,
colorlinks=true,
linkcolor=red,
citecolor=blue,
filecolor=magenta,
urlcolor=cyan
}

\usepackage{array}

\externaldocument{../../SRS/SRS}

\input{../../Comments.text}
\input{../../Common.text}

\begin{document}

\title{Module Interface Specification for \progname{}}

\author{\authname}

\date{\today}

\maketitle

\pagenumbering{roman}

\section{Revision History}

\begin{tabularx}{\textwidth}{p{3cm}p{2cm}X}
\toprule {\bf Date} & {\bf Version} & {\bf Notes}\\
\midrule
2026-03-23 & 1.0 & Initial MG/MIS presentation draft\\
2026-04-02 & 1.1 & Full MIS formalization aligned with 8-module implementation\\
\bottomrule
\end{tabularx}

~\newpage

\section{Symbols, Abbreviations and Acronyms}

\renewcommand{\arraystretch}{1.2}
\begin{tabular}{l l}
\toprule
\textbf{Symbol} & \textbf{Description}\\
\midrule
MIS & Module Interface Specification \\
MG  & Module Guide \\
SRS & Software Requirements Specification \\
ABC & Abstract Base Class (Python interface enforcement) \\
M1--M8 & Module identifiers as defined in the MG \\
$\mathbb{R}$ & Real number \\
$\mathbb{Z}$ & Integer \\
$\mathbb{N}$ & Natural number \\
Point2D & A tuple $(x, y) \in \mathbb{R}^2$ representing image coordinates \\
Point3D & A tuple $(x, y, z) \in \mathbb{R}^3$ representing 3D coordinates \\
JointID & An integer index identifying a MediaPipe pose landmark \\
LandmarkList & A sequence of 33 Point3D landmarks \\
\bottomrule
\end{tabular}

\newpage

\tableofcontents

\newpage

\pagenumbering{arabic}

\section{Introduction}

The following document details the Module Interface Specifications for
\progname{}, an augmented reality application that provides real-time
biomechanical feedback for fitness exercise using a standard webcam.

Complementary documents include the System Requirement Specifications
and Module Guide. The full documentation and implementation can be
found at \url{https://github.com/mmdmirh/cas741}.

\section{Notation}

The structure of the MIS for modules comes from \citet{HoffmanAndStrooper1995},
with the addition that template modules have been adapted from
\cite{GhezziEtAl2003}. The mathematical notation comes from Chapter 3 of
\citet{HoffmanAndStrooper1995}. The symbol := is used for assignment and
conditional rules follow the form $(c_1 \Rightarrow r_1 \mid c_2 \Rightarrow r_2)$.

The following table summarizes the primitive data types used by \progname.

\begin{center}
\renewcommand{\arraystretch}{1.2}
\noindent
\begin{tabular}{l l p{7.5cm}}
\toprule
\textbf{Data Type} & \textbf{Notation} & \textbf{Description}\\
\midrule
character & char & a single symbol or digit\\
integer & $\mathbb{Z}$ & a number without a fractional component in $(-\infty, \infty)$ \\
natural number & $\mathbb{N}$ & a number without a fractional component in $[1, \infty)$ \\
real & $\mathbb{R}$ & any number in $(-\infty, \infty)$\\
boolean & $\mathbb{B}$ & \texttt{True} or \texttt{False}\\
\bottomrule
\end{tabular}
\end{center}

\noindent
Derived types used include: sequences (ordered lists of same-type elements),
tuples (fixed-length heterogeneous collections), and functions (typed
input/output mappings). Interfaces are enforced in Python via Abstract Base
Classes (ABCs) from the \texttt{abc} module.

\section{Module Decomposition}

The following table is taken directly from the Module Guide document for this project.

\begin{table}[h!]
\centering
\begin{tabular}{p{0.3\textwidth} p{0.6\textwidth}}
\toprule
\textbf{Level 1} & \textbf{Level 2}\\
\midrule

{Hardware-Hiding} & M1: Video Input Module \\
\midrule

\multirow{4}{0.3\textwidth}{Behaviour-Hiding} & M2: Display Output Module\\
& M3: Video Formatting Module\\
& M6: Exercise State Module\\
& M7: UI Rendering Module\\
\midrule

\multirow{3}{0.3\textwidth}{Software Decision} & M4: Pose Tracking Module\\
& M5: Kinematic Engine Module\\
& M8: Signal Smoothing Module\\
\bottomrule

\end{tabular}
\caption{Module Hierarchy}
\label{TblMH}
\end{table}

\newpage

%-----------------------------------------------------------------------
\section{MIS of M3: Video Formatting Module} \label{M3}
%-----------------------------------------------------------------------

\subsection{Module}
\texttt{m3\_video\_formatting} (\texttt{backend/modules/m3\_video\_formatting.py})

\subsection{Uses}
M1 (Video Input)

\subsection{Syntax}

\subsubsection{Exported Constants}
None.

\subsubsection{Exported Access Programs}

\begin{center}
\begin{tabular}{p{3.5cm} p{4cm} p{3cm} p{2cm}}
\hline
\textbf{Name} & \textbf{In} & \textbf{Out} & \textbf{Exceptions} \\
\hline
\texttt{encode\_frame} & frame: bytes & encoded: str & \texttt{ValueError} \\
\texttt{decode\_frame} & encoded: str & frame: bytes & \texttt{ValueError} \\
\hline
\end{tabular}
\end{center}

\subsection{Semantics}

\subsubsection{State Variables}
None. (Stateless conversion module.)

\subsubsection{Environment Variables}
Raw JPEG/PNG image data from the webcam feed.

\subsubsection{Assumptions}
Input bytes represent a valid compressed image (JPEG or PNG).

\subsubsection{Access Routine Semantics}

\noindent \texttt{encode\_frame(frame)}:
\begin{itemize}
  \item output: Base64-encoded UTF-8 string representation of the raw image bytes.
  \item exception: \texttt{ValueError} if \texttt{frame} is empty or not valid image bytes.
\end{itemize}

\noindent \texttt{decode\_frame(encoded)}:
\begin{itemize}
  \item output: Raw image bytes decoded from the Base64 string.
  \item exception: \texttt{ValueError} if \texttt{encoded} is not valid Base64.
\end{itemize}

\newpage

%-----------------------------------------------------------------------
\section{MIS of M4: Pose Tracking Module} \label{M4}
%-----------------------------------------------------------------------

\subsection{Module}
\texttt{m4\_pose\_tracking} (\texttt{backend/modules/m4\_pose\_tracking.py})

\subsection{Uses}
M3 (Video Formatting), M8 (Signal Smoothing)

\subsection{Syntax}

\subsubsection{Exported Constants}
\texttt{NUM\_LANDMARKS} := 33 (number of MediaPipe Pose landmarks)

\subsubsection{Exported Access Programs}

\begin{center}
\begin{tabular}{p{3.5cm} p{4cm} p{3.5cm} p{2cm}}
\hline
\textbf{Name} & \textbf{In} & \textbf{Out} & \textbf{Exceptions} \\
\hline
\texttt{detect} & frame: bytes & LandmarkList & \texttt{DetectionError}\\
\texttt{is\_visible} & landmarks: LandmarkList, id: JointID & $\mathbb{B}$ & - \\
\hline
\end{tabular}
\end{center}

\subsection{Semantics}

\subsubsection{State Variables}
\begin{itemize}
  \item \texttt{model}: Loaded ML inference model (MediaPipe Pose instance).
\end{itemize}

\subsubsection{Environment Variables}
\begin{itemize}
  \item Webcam image frame (from M3).
\end{itemize}

\subsubsection{Assumptions}
\begin{itemize}
  \item The input frame contains at most one person.
  \item Lighting conditions meet minimum MediaPipe requirements.
\end{itemize}

\subsubsection{Access Routine Semantics}

\noindent \texttt{detect(frame)}:
\begin{itemize}
  \item output: A sequence of 33 landmarks, each with $(x, y, z, \text{visibility})$ where $x, y, z \in [0.0, 1.0]$ normalized to image dimensions.
  \item exception: \texttt{DetectionError} if the frame cannot be decoded or no person is detected.
\end{itemize}

\noindent \texttt{is\_visible(landmarks, id)}:
\begin{itemize}
  \item output: \texttt{True} iff \texttt{landmarks[id].visibility} $> 0.5$.
\end{itemize}

\newpage

%-----------------------------------------------------------------------
\section{MIS of M5: Kinematic Engine Module} \label{M5}
%-----------------------------------------------------------------------

\subsection{Module}
\texttt{m5\_kinematic\_engine} (\texttt{backend/modules/m5\_kinematic\_engine.py})

\subsection{Uses}
M4 (Pose Tracking)

\subsection{Syntax}

\subsubsection{Exported Constants}
None.

\subsubsection{Exported Access Programs}

\begin{center}
\begin{tabular}{p{4cm} p{4.5cm} p{2.5cm} p{2.5cm}}
\hline
\textbf{Name} & \textbf{In} & \textbf{Out} & \textbf{Exceptions} \\
\hline
\texttt{compute\_angle} & a, b, c: Point3D & $\mathbb{R}$ & \texttt{ValueError}\\
\texttt{extract\_metrics} & landmarks: LandmarkList, exercise: str & dict & \texttt{KeyError}\\
\hline
\end{tabular}
\end{center}

\subsection{Semantics}

\subsubsection{State Variables}
None. (Stateless computation module.)

\subsubsection{Environment Variables}
None.

\subsubsection{Assumptions}
\begin{itemize}
  \item All three points $a$, $b$, $c$ are non-collinear (i.e., $\|a - b\| > 0$ and $\|c - b\| > 0$).
  \item \texttt{exercise} string is one of the recognized exercise types (e.g., \texttt{"squat"}, \texttt{"curl"}).
\end{itemize}

\subsubsection{Access Routine Semantics}

\noindent \texttt{compute\_angle(a, b, c)}:
\begin{itemize}
  \item output: The angle $\theta$ at joint $b$ formed by vectors $\overrightarrow{ba}$ and $\overrightarrow{bc}$, in degrees:
  $$\theta = \arccos\left(\frac{\overrightarrow{ba} \cdot \overrightarrow{bc}}{\|\overrightarrow{ba}\| \cdot \|\overrightarrow{bc}\|}\right) \times \frac{180}{\pi}$$
  where $\theta \in [0, 180]$.
  \item exception: \texttt{ValueError} if any input vector has zero magnitude.
\end{itemize}

\noindent \texttt{extract\_metrics(landmarks, exercise)}:
\begin{itemize}
  \item output: A dictionary \texttt{\{"primary": $\mathbb{R}$, "form": dict\}} where:
  \begin{itemize}
    \item \texttt{primary}: The main joint angle driving the rep count.
    \item \texttt{form}: A dictionary of secondary form metrics (e.g., \texttt{knee\_valgus\_index}, \texttt{torso\_lean}).
  \end{itemize}
  \item exception: \texttt{KeyError} if \texttt{exercise} is not a recognized type.
\end{itemize}

\subsubsection{Local Functions}

\noindent \texttt{\_get\_coords(landmarks, id: JointID) $\rightarrow$ Point3D}:
\begin{itemize}
  \item Returns $(x, y, z)$ for landmark at index \texttt{id}.
\end{itemize}

\noindent \texttt{\_compute\_distance(p1, p2: Point3D) $\rightarrow \mathbb{R}$}:
\begin{itemize}
  \item Returns Euclidean distance $\|p1 - p2\|$.
\end{itemize}

\newpage

%-----------------------------------------------------------------------
\section{MIS of M6: Exercise State Module} \label{M6}
%-----------------------------------------------------------------------

\subsection{Module}
\texttt{m6\_exercise\_state} (\texttt{backend/modules/m6\_exercise\_state.py})

\subsection{Uses}
M5 (Kinematic Engine)

\subsection{Syntax}

\subsubsection{Exported Constants}
\begin{itemize}
  \item \texttt{STATES} := \{\texttt{BOTTOM}, \texttt{UP\_PHASE}, \texttt{TOP}, \texttt{DOWN\_PHASE}\}
\end{itemize}

\subsubsection{Exported Access Programs}

\begin{center}
\begin{tabular}{p{3.5cm} p{4cm} p{3cm} p{2cm}}
\hline
\textbf{Name} & \textbf{In} & \textbf{Out} & \textbf{Exceptions} \\
\hline
\texttt{update} & angle: $\mathbb{R}$, dt: $\mathbb{R}$ & state: str & - \\
\texttt{get\_rep\_count} & - & $\mathbb{N}$ & - \\
\texttt{reset} & - & - & - \\
\hline
\end{tabular}
\end{center}

\subsection{Semantics}

\subsubsection{State Variables}
\begin{itemize}
  \item \texttt{current\_state}: str $\in$ \texttt{STATES}
  \item \texttt{rep\_count}: $\mathbb{N}$
  \item \texttt{angle\_min}: $\mathbb{R}$ (minimum angle observed in current rep)
  \item \texttt{angle\_max}: $\mathbb{R}$ (maximum angle observed in current rep)
\end{itemize}

\subsubsection{Environment Variables}
None.

\subsubsection{Assumptions}
\begin{itemize}
  \item Angles are produced by M5 at a rate sufficient for state detection ($\geq$ 15 fps).
  \item Calibration parameters (ROM thresholds) are set before \texttt{update} is called.
\end{itemize}

\subsubsection{Access Routine Semantics}

\noindent \texttt{update(angle, dt)}:
\begin{itemize}
  \item transition: Advances the state machine according to:
  \begin{align*}
    \texttt{BOTTOM} &\xrightarrow{\text{angle} \geq \theta_{top\_min}} \texttt{UP\_PHASE} \\
    \texttt{UP\_PHASE} &\xrightarrow{\text{angle} \geq \theta_{top\_max}} \texttt{TOP} \\
    \texttt{TOP} &\xrightarrow{\text{angle} \leq \theta_{bottom\_max}} \texttt{DOWN\_PHASE} \\
    \texttt{DOWN\_PHASE} &\xrightarrow{\text{angle} \leq \theta_{bottom\_min}} \texttt{BOTTOM},\; \texttt{rep\_count} := \texttt{rep\_count} + 1
  \end{align*}
  \item output: The new \texttt{current\_state} string after the transition.
\end{itemize}

\noindent \texttt{get\_rep\_count()}:
\begin{itemize}
  \item output: Current value of \texttt{rep\_count}.
\end{itemize}

\noindent \texttt{reset()}:
\begin{itemize}
  \item transition: \texttt{rep\_count} := 0, \texttt{current\_state} := \texttt{BOTTOM}.
\end{itemize}

\newpage

%-----------------------------------------------------------------------
\section{MIS of M8: Signal Smoothing Module} \label{M8}
%-----------------------------------------------------------------------

\subsection{Module}
\texttt{m8\_signal\_smoothing} (\texttt{backend/modules/m8\_signal\_smoothing.py})

\subsection{Uses}
M4 (Pose Tracking)

\subsection{Syntax}

\subsubsection{Exported Constants}
None.

\subsubsection{Exported Access Programs}

\begin{center}
\begin{tabular}{p{3.5cm} p{4.5cm} p{3cm} p{2cm}}
\hline
\textbf{Name} & \textbf{In} & \textbf{Out} & \textbf{Exceptions} \\
\hline
\texttt{smooth} & measurement: $\mathbb{R}$ & estimate: $\mathbb{R}$ & - \\
\texttt{smooth\_landmarks} & landmarks: LandmarkList & LandmarkList & - \\
\hline
\end{tabular}
\end{center}

\subsection{Semantics}

\subsubsection{State Variables}
\begin{itemize}
  \item \texttt{x\_est}: $\mathbb{R}$ --- current state estimate
  \item \texttt{P}: $\mathbb{R}$ --- estimation error covariance
  \item \texttt{Q}: $\mathbb{R}$ --- process noise covariance
  \item \texttt{R}: $\mathbb{R}$ --- measurement noise covariance
\end{itemize}

\subsubsection{Environment Variables}
None.

\subsubsection{Assumptions}
\begin{itemize}
  \item Measurements are scalar real values from M4 at a fixed frame rate.
  \item Process and measurement noise covariances $Q$ and $R$ are tuned to match webcam noise characteristics.
\end{itemize}

\subsubsection{Access Routine Semantics}

\noindent \texttt{smooth(measurement)}:
\begin{itemize}
  \item transition: One-step Kalman filter update:
  \begin{align*}
    K &:= P / (P + R) & \text{(Kalman Gain)}\\
    x\_est &:= x\_est + K \times (measurement - x\_est) & \text{(State Update)}\\
    P &:= (1 - K) \times P + Q & \text{(Covariance Update)}
  \end{align*}
  \item output: Updated \texttt{x\_est} (the filtered scalar estimate).
\end{itemize}

\noindent \texttt{smooth\_landmarks(landmarks)}:
\begin{itemize}
  \item output: A new LandmarkList where each $(x, y, z)$ coordinate of every landmark has been individually filtered by \texttt{smooth()}.
\end{itemize}

\newpage

\bibliographystyle {plainnat}
\bibliography {../../../refs/References}

\end{document}
